\begin{proof}
    Let $( X, Y ) \in  Q$.
    There are 9 cases to consider.
    \begin{enumerate}
    \item Assume that $( X, Y )$ is rule Q(1).
          Clearly $\scalebox{0.65}{\circimg{tikz-cache/hqc__eqn7_lhs.pdf}} \in \decoder( X )$ and $0 \in \decoder( Y )$.
          It then follows by $\HqcPOrd$ that $\scalebox{0.65}{\circimg{tikz-cache/hqc__eqn7_lhs.pdf}} \approx_E 0$.
    \item Assume that $( X, Y )$ is rule Q(2).
          Clearly $\scalebox{0.65}{\circimg{tikz-cache/proofs__q2_1.pdf}} \in \decoder( X )$ and $\scalebox{0.65}{\circimg{tikz-cache/proofs__q2_3.pdf}} \in \decoder( Y )$.
          It then follows by $\HqcPowAdd$ that $\scalebox{0.65}{\circimg{tikz-cache/proofs__q2_1.pdf}} \approx_E \scalebox{0.65}{\circimg{tikz-cache/proofs__q2_2.pdf}} = \scalebox{0.65}{\circimg{tikz-cache/proofs__q2_3.pdf}}$.
    \item Assume that $( X, Y )$ is rule Q(3).
          Clearly $\sigma \in \decoder( X )$.
          Since $\CatCQC$ is a controlled prop category, then the following equation holds.
          \begin{align*}
              Y
              =
              \scalebox{0.65}{\circimg{tikz-cache/cqc__eqn3_rhs.pdf}}
              &=
              \scalebox{0.65}{\circimg{tikz-cache/proofs__simp_eqn3_1.pdf}}
              \\
              &=
              \scalebox{0.65}{\circimg{tikz-cache/proofs__simp_eqn3_2.pdf}}
          \end{align*}
          Then by \cref{Lemma:CNot}, the following equation also holds.
          \begin{equation*}
              \scalebox{0.65}{\circimg{tikz-cache/proofs__decode_q3_rhs.pdf}}
              \in
              \decoder( Y )
          \end{equation*}
          Moreover, the following derivation holds.
          \begin{equation*}
              \scalebox{0.65}{\circimg{tikz-cache/proofs__decode_q3_lhs.pdf}}
              \overset{\HqcSwapDef}{\approx}
              \scalebox{0.65}{\circimg{tikz-cache/proofs__decode_q3_mid.pdf}}
              \overset{\HqcCtrlSwap}{\approx}
              \scalebox{0.65}{\circimg{tikz-cache/proofs__decode_q3_rhs.pdf}}
          \end{equation*}
    \item Assume that $( X, Y )$ is rule Q(4).
          Clearly $\scalebox{0.65}{\circimg{tikz-cache/cqc__eqn4_lhs.pdf}} \in \decoder( X )$ and $\scalebox{0.65}{\circimg{tikz-cache/cqc__eqn4_rhs.pdf}} \in \decoder( Y )$.
          Since $H = H^\dagger$, then $\scalebox{0.65}{\circimg{tikz-cache/cqc__eqn4_lhs.pdf}} \approx_E \scalebox{0.65}{\circimg{tikz-cache/cqc__eqn4_rhs.pdf}}$ by \cref{Thm:HQCInv}.
    \item Assume that $( X, Y )$ is rule Q(5).
          Then $\alpha_1 \in \mathbb{R}$, $\alpha_2 \in \mathbb{R}$, and $( \beta_0, \beta_1, \beta_2, \beta_2 ) = \mathsf{Euler}( \alpha_1, \alpha_2 )$.
          Clearly the following equations hold (see~\cref{App:CNOT}).
          \begin{align*}
              \scalebox{0.65}{\circimg{tikz-cache/proofs__euler_1.pdf}}
              &\in
              \decoder\left( \scalebox{0.65}{\circimg{tikz-cache/cqc__eqn5_lhs.pdf}} \right)
              \\
              \scalebox{0.65}{\circimg{tikz-cache/proofs__euler_11.pdf}}
              &\in
              \decoder\left( \scalebox{0.65}{\circimg{tikz-cache/cqc__eqn5_rhs.pdf}} \right)
          \end{align*}
          Then the following derivation holds since $H = H^\dagger$.
          \begin{align*}
              \scalebox{0.65}{\circimg{tikz-cache/proofs__euler_1.pdf}}
              \overset{\HqcFreePhase}{\approx}&
              \scalebox{0.65}{\circimg{tikz-cache/proofs__euler_2.pdf}}
              \\
              \overset{\HqcFreePhase}{\approx}&
              \scalebox{0.65}{\circimg{tikz-cache/proofs__euler_3.pdf}}
              \\
              \overset{\HqcPowDag}{\approx}&
              \scalebox{0.65}{\circimg{tikz-cache/proofs__euler_4.pdf}}
              \\
              \overunderset{\HqcPowCong}{\HqcTargMDef}{\approx}&
              \scalebox{0.65}{\circimg{tikz-cache/proofs__euler_5.pdf}}
              \\
              \overset{\HqcEuler}{\approx}\,&
              \scalebox{0.65}{\circimg{tikz-cache/proofs__euler_6.pdf}}
              \\
              \overset{\HqcFreePhase}{\approx}&
              \scalebox{0.65}{\circimg{tikz-cache/proofs__euler_7.pdf}}
              \\
              \overunderset{\HqcPowCong}{\HqcTargMDef}{\approx}&
              \scalebox{0.65}{\circimg{tikz-cache/proofs__euler_8.pdf}}
              \\
              \overset{\HqcPowDag}{\approx}&
              \scalebox{0.65}{\circimg{tikz-cache/proofs__euler_9.pdf}}
              \\
              \overset{\HqcFreePhase}{\approx}&
              \scalebox{0.65}{\circimg{tikz-cache/proofs__euler_10.pdf}}
              \\
              \overset{\HqcFreePhase}{\approx}&
              \scalebox{0.65}{\circimg{tikz-cache/proofs__euler_11.pdf}}
          \end{align*}
    \item Assume that $( X, Y )$ is rule Q(6).
          As in~\ref{App:CNOT}, the following equations hold where $\beta = \alpha / \pi$.
          \begin{align*}
              \scalebox{0.65}{\circimg{tikz-cache/proofs__q6_1.pdf}}
              &\in
              \decoder\left( \scalebox{0.65}{\circimg{tikz-cache/cqc__eqn6_lhs.pdf}} \right)
              &
              \scalebox{0.65}{\circimg{tikz-cache/proofs__q6_4.pdf}}
              &\in
              \decoder\left( \scalebox{0.65}{\circimg{tikz-cache/cqc__eqn6_rhs.pdf}} \right)
          \end{align*}
          Since $C( - )$ is a conjugated control functor and $( H \boxtimes 0 )^\dagger = H^\dagger \boxtimes 0^\dagger = H \boxtimes 0$, then the following derivation holds.
          \begin{equation*}
              \scalebox{0.65}{\circimg{tikz-cache/proofs__q6_1.pdf}}
              \approx_E
              \scalebox{0.65}{\circimg{tikz-cache/proofs__q6_2.pdf}}
              \approx_E
              \scalebox{0.65}{\circimg{tikz-cache/proofs__q6_3.pdf}}
              \approx_E
              \scalebox{0.65}{\circimg{tikz-cache/proofs__q6_4.pdf}}
          \end{equation*}
    \item Assume that $( X, Y )$ is rule Q(7).
          As in~\ref{App:CNOT}, the following equations hold where $\beta = \alpha / \pi$.
          \begin{align*}
              \scalebox{0.65}{\circimg{tikz-cache/proofs__q7_1.pdf}}
              &\in
              \decoder\left( \scalebox{0.65}{\circimg{tikz-cache/cqc__eqn7_lhs.pdf}} \right)
              &
              \scalebox{0.65}{\circimg{tikz-cache/proofs__q7_5.pdf}}
              &\in
              \decoder\left( \scalebox{0.65}{\circimg{tikz-cache/cqc__eqn7_rhs.pdf}} \right)
          \end{align*}
          Since $C( - )$ is a conjugated control functor and $( U \exparr \alpha )^\dagger = U \exparr ( -\alpha )$ for each unitary $U$, then the following derivation holds.
          \begin{align*}
              \scalebox{0.65}{\circimg{tikz-cache/proofs__q7_1.pdf}}
              \overunderset{\HQCCtrlCong}{\HqcFreePhase}{\approx}
              \scalebox{0.61}{\circimg{tikz-cache/proofs__q7_2.pdf}}
              \overunderset{\HQCCtrlCong}{\HqcFreePhase}{\approx}
              \scalebox{0.65}{\circimg{tikz-cache/proofs__q7_3.pdf}}
              \approx_E
              \scalebox{0.65}{\circimg{tikz-cache/proofs__q7_4.pdf}}
              \approx_E
              \scalebox{0.65}{\circimg{tikz-cache/proofs__q7_5.pdf}}
          \end{align*}
    \item Assume that $( X, Y )$ is rule Q(8).
          Clearly the following equations hold (see~\cref{App:CNOT}).
          \begin{align*}
              \scalebox{0.65}{\circimg{tikz-cache/proofs__q8_1.pdf}}
              &\in
              \decoder\left( \scalebox{0.65}{\circimg{tikz-cache/cqc__eqn8_lhs.pdf}} \right)
              &
              \scalebox{0.65}{\circimg{tikz-cache/proofs__q8_4.pdf}}
              &\in
              \decoder\left( \scalebox{0.65}{\circimg{tikz-cache/cqc__eqn8_rhs.pdf}} \right)
          \end{align*}
          Since $C( - )$ is a conjugated control functor and $( H \boxtimes 1 \boxtimes 0 )^\dagger = H^\dagger \boxtimes 1^\dagger \boxtimes 0^\dagger = H \boxtimes 1 \boxtimes 0$, then the following derivation holds.
          \begin{equation*}
              \scalebox{0.65}{\circimg{tikz-cache/proofs__q8_1.pdf}}
              \approx_E
              \scalebox{0.65}{\circimg{tikz-cache/proofs__q8_2.pdf}}
              \approx_E
              \scalebox{0.65}{\circimg{tikz-cache/proofs__q8_3.pdf}}
              \approx_E
              \scalebox{0.65}{\circimg{tikz-cache/proofs__q8_4.pdf}}
          \end{equation*}
    \item Assume that $( X, Y )$ is rule Q(9).
          Clearly the following equations hold (see~\cref{App:CNOT}).
          \begin{align*}
              \scalebox{0.65}{\circimg{tikz-cache/proofs__q9_1.pdf}}
              &\in
              \decoder\left( \scalebox{0.65}{\circimg{tikz-cache/cqc__eqn9_lhs.pdf}} \right)
              &
              \scalebox{0.65}{\circimg{tikz-cache/proofs__q9_8.pdf}}
              &\in
              \decoder\left( \scalebox{0.65}{\circimg{tikz-cache/cqc__eqn9_rhs.pdf}} \right)
          \end{align*}
          As in the previous proofs, this case will be handled using the property that $C( - )$ is a conjugated control functor.
          First, it must be shown that the third gate on the left-hand side is the adjoint to the first gate on the left-hand side.
          \begin{align*}
              \scalebox{0.65}{\circimg{tikz-cache/proofs__q9_adj_1.pdf}}
              \overunderset{\HQCCtrlCong}{\HqcFreePhase}{\approx}&
              \scalebox{0.65}{\circimg{tikz-cache/proofs__q9_adj_2.pdf}}
              \overset{\text{Thm(\ref{Thm:HQCInv})}}{\approx}
              \scalebox{0.65}{\circimg{tikz-cache/proofs__q9_adj_3.pdf}}
              \overset{\HqcCtrlComp}{\approx}
              \scalebox{0.65}{\circimg{tikz-cache/proofs__q9_adj_4.pdf}}
              \\
              \overunderset{\HQCCtrlCong}{\HqcPowAdd}{\approx}&
              \scalebox{0.65}{\circimg{tikz-cache/proofs__q9_adj_5.pdf}}
              \overunderset{\HQCCtrlCong}{\HqcCtrlId}{\approx}
              \scalebox{0.65}{\circimg{tikz-cache/proofs__q9_adj_6.pdf}}
              \overset{\HqcCtrlId}{\approx}
              \scalebox{0.65}{\circimg{tikz-cache/proofs__q9_adj_7.pdf}}
          \end{align*}
          Let $U$ denote the final string in this derivation.
          Since $C( - )$ is a conjugated control functor and $C( V )^\dagger = C( V ) \exparr ( -1 )$ for each unitary $V$, then the following derivation holds.
          \begin{align*}
              \scalebox{0.65}{\circimg{tikz-cache/proofs__q9_1.pdf}}
              \overunderset{\HQCCtrlCong}{\HqcFreePhase}{\approx}&
              \scalebox{0.65}{\circimg{tikz-cache/proofs__q9_2.pdf}}
              \approx_E
              \scalebox{0.65}{\circimg{tikz-cache/proofs__q9_3.pdf}}
              \approx_E
              \scalebox{0.65}{\circimg{tikz-cache/proofs__q9_4.pdf}}
              \approx_E
              \scalebox{0.65}{\circimg{tikz-cache/proofs__q9_5.pdf}}
              \\
              \approx_E&
              \scalebox{0.65}{\circimg{tikz-cache/proofs__q9_6.pdf}}
              \overunderset{\HQCCtrlCong}{\HqcFreePhase}{\approx}
              \scalebox{0.65}{\circimg{tikz-cache/proofs__q9_7.pdf}}
              \overunderset{\HQCCtrlCong}{\HqcFreePhase}{\approx}
              \scalebox{0.65}{\circimg{tikz-cache/proofs__q9_8.pdf}}
          \end{align*}
    \end{enumerate}
    In each case, there exists a $U \in \decoder( X )$ and $V \in \decoder( Y )$ such that $U \approx_E V$.
    Since $( X, Y )$ was arbitrary, then this completes the proof.
\end{proof}
